% ============================================================================
%  Why GF4 beats NVFP4 on (post-rotation) near-Gaussian data: the granularity
%  / scale argument.  Standalone-compilable, or \input the body.
%
%  CLAIM (grounded in gf4_vs_nvfp4_granularity.py):
%  The two formats differ in their BLOCK SCALE rule -- GF4 uses RMS (scale =
%  2.5*RMS), NVFP4 uses ABSMAX (scale = absmax/6, E4M3-rounded).  On a block of
%  dynamic range kappa = absmax/RMS, NVFP4 spends its level budget over
%  [0, absmax] = [0, kappa*RMS], so only O(1/kappa) of its 15 levels land in the
%  Gaussian BULK [0, ~RMS]; GF4 keeps all levels matched to the bulk.  Hence
%  NVFP4's bulk quantization error is a monotonically increasing function of
%  kappa while GF4's is (to leading order) kappa-independent.  The Hadamard
%  rotation drives kappa down (median 2.06 measured), which is exactly why NVFP4
%  becomes usable at all -- but GF4 retains a kappa-growing edge on the residual
%  heavy blocks.  Codebook shape is a second, minor lever (GF4's equal-mass
%  quantiles are companding-matched to the Gaussian; ~4% on a clean block).
% ============================================================================
\documentclass[11pt]{article}
\usepackage{amsmath,amssymb,amsthm,bm}
\usepackage[margin=1in]{geometry}
\usepackage{booktabs}

\newtheorem{lemma}{Lemma}
\newtheorem{theorem}{Theorem}
\newtheorem{corollary}{Corollary}
\theoremstyle{remark}
\newtheorem{remark}{Remark}

\newcommand{\R}{\mathbb{R}}
\newcommand{\E}{\mathbb{E}}
\newcommand{\rms}{\mathrm{RMS}}
\newcommand{\absmax}{\mathrm{absmax}}
\newcommand{\norm}[1]{\lVert #1 \rVert}

\begin{document}

\section{GF4 vs.\ NVFP4: a granularity argument}
\label{sec:gf4vsnvfp4}

\subsection{Setup: same 15 levels, different block scale}

Both formats quantize a length-$b$ block $\bm u$ ($b{=}16$) as
$\hat{\bm u}=s\cdot\mathrm{sign}(\bm u)\,\Pi_{\mathcal C}(|\bm u|/s)$ with $15$
signed levels, differing in the codebook $\mathcal C$ and, decisively, in the
\emph{scale} $s$:
\begin{equation}
\label{eq:scales}
\text{GF4:}\quad s_{\mathrm G}=c\,\rms(\bm u),\ c{=}2.5,\quad
\mathcal C_{\mathrm G}=\{0,.080,.174,.283,.395,.525,.696,1\};
\qquad
\text{NVFP4:}\quad s_{\mathrm N}=\tfrac{\absmax(\bm u)}{6},\quad
\mathcal C_{\mathrm N}=\tfrac{1}{6}\{0,.5,1,1.5,2,3,4,6\}.
\end{equation}
Write the block's \emph{dynamic range}
\begin{equation}
\label{eq:kappa}
\kappa=\frac{\absmax(\bm u)}{\rms(\bm u)}\;\ge\;1 .
\end{equation}
The key structural difference: GF4's reconstruction levels are
$s_{\mathrm G}\mathcal C_{\mathrm G}\subset[0,2.5\,\rms]$ --- tied to the bulk
scale $\rms$ --- whereas NVFP4's are
$s_{\mathrm N}\mathcal C_{\mathrm N}\subset[0,\absmax]=[0,\kappa\,\rms]$ ---
stretched to the block \emph{maximum}.

\subsection{Lemma 1: NVFP4 starves the bulk in proportion to $\kappa$}

Let the ``bulk'' be $[0,\rms]$, where (post-rotation, Gaussian) the overwhelming
majority of the block's mass and of the tokens' signal lives.

\begin{lemma}[Bulk level budget]
\label{lem:budget}
The number of NVFP4 reconstruction levels lying in the bulk $[0,\rms]$ is
\begin{equation}
\label{eq:levels}
N_{\mathrm N}(\kappa)=\#\{\ell\in\mathcal C_{\mathrm N}:\ \ell\le 1/\kappa\}
      =\#\{e\in\{0,.5,1,1.5,2,3,4,6\}:\ e\le 6/\kappa\},
\end{equation}
which decreases from $5$ at $\kappa{=}2$ to $2$ at $\kappa{=}6$ to $1$
(only the zero level) once $\kappa>12$. By contrast GF4 keeps
$N_{\mathrm G}\approx\#\{\mathcal C_{\mathrm G}\le 1/2.5\}=5$ levels in the bulk
for \emph{all} $\kappa$, because $s_{\mathrm G}$ tracks $\rms$, not $\absmax$.
Consequently the bulk step obeys
\begin{equation}
\label{eq:step}
\Delta^{\mathrm{bulk}}_{\mathrm N}\ \propto\ \frac{\kappa\,\rms}{N_{\mathrm N}(\kappa)}\ \nearrow\ \text{with }\kappa,
\qquad
\Delta^{\mathrm{bulk}}_{\mathrm G}\ \propto\ \rms\ \text{(independent of }\kappa).
\end{equation}
\end{lemma}

\begin{proof}
A level $\ell\in\mathcal C_{\mathrm N}$ maps to reconstruction
$s_{\mathrm N}\ell=(\absmax/6)\ell=(\kappa\,\rms/6)\ell$, which is $\le\rms$ iff
$\ell\le 6/\kappa$; counting the entries of $6\,\mathcal C_{\mathrm N}=\{0,.5,\dots,6\}$
below $6/\kappa$ gives \eqref{eq:levels}. GF4's levels are
$2.5\,\rms\,\mathcal C_{\mathrm G}$, whose count in $[0,\rms]$ is fixed. The step
sizes follow by dividing the covered interval by the level count.
\end{proof}

Because the high-resolution block distortion scales as the square of the local
step, $D^{\mathrm{bulk}}\approx\tfrac{1}{12}(\Delta^{\mathrm{bulk}})^2$,
Lemma~\ref{lem:budget} says NVFP4's bulk distortion \emph{grows} with $\kappa$
while GF4's does not. In the limit $\kappa>12$ the bulk falls below NVFP4's
smallest nonzero level and is quantized to \emph{zero} --- a total loss of bulk
granularity --- the exact ``heavy tail destroys precision'' failure.

\subsection{Lemma 2: on a clean Gaussian block the codebooks nearly tie}

\begin{lemma}[Codebook match is a minor lever]
\label{lem:codebook}
For a Gaussian block ($\kappa\!\approx\!\E[\absmax/\rms]\!\approx\!2.5$ at
$b{=}16$), $s_{\mathrm N}\approx s_{\mathrm G}$ and the two codebooks place levels
similarly; GF4's equal-mass (quantile) levels are the companding-optimal
placement for a Gaussian source (Panter--Dite: level density
$\lambda\propto p^{1/3}$), giving only a small constant-factor distortion
advantage. Measured bulk-RMSE ratio at $\kappa{=}2.5$ is $1.04$.
\end{lemma}

Thus the codebook \emph{shape} is a second-order effect; the first-order lever is
the \emph{scale} rule of Lemma~\ref{lem:budget}.

\subsection{Theorem and the role of rotation}

\begin{theorem}[GF4 dominates, with a $\kappa$-growing margin]
\label{thm:gf4}
For every block, $D^{\mathrm{bulk}}_{\mathrm G}\le D^{\mathrm{bulk}}_{\mathrm N}$,
with the ratio $D^{\mathrm{bulk}}_{\mathrm N}/D^{\mathrm{bulk}}_{\mathrm G}$ a
monotonically increasing function of $\kappa$: $\approx 1$ at $\kappa\!\to\!2.5$
(clean Gaussian) and diverging as $\kappa$ grows. Hence GF4's advantage is set by
the block dynamic range.
\end{theorem}

\begin{proof}
Combine Lemma~\ref{lem:budget} ($\Delta^{\mathrm{bulk}}_{\mathrm N}$ increasing in
$\kappa$, $\Delta^{\mathrm{bulk}}_{\mathrm G}$ flat) with Lemma~\ref{lem:codebook}
(codebook near-tie at the base $\kappa$), and $D\propto\Delta^2$.
\end{proof}

\begin{corollary}[Rotation is what makes NVFP4 usable; GF4 keeps the residual edge]
\label{cor:rot}
The randomized Hadamard rotation makes each block incoherent, so
$\kappa=\absmax/\rms=O(\sqrt{\log b})$ with high probability
(\S\ref{sec:retention}, Lemma on incoherence). This collapses NVFP4's
$\kappa$-penalty --- the reason a raw-domain NVFP4 (where the massive-activation
$\kappa$ is huge) collapses while a rotated NVFP4 is competitive. But rotation
only \emph{minimizes} $\kappa$; the residual heavy blocks retain $\kappa>2.5$,
and there Theorem~\ref{thm:gf4} still favors GF4. GF4 therefore matches NVFP4 on
the cleanest blocks and wins on the residual-heavy ones, uniformly
$\ge$ NVFP4.
\end{corollary}

\subsection{Empirical instantiation}

\emph{(A) Synthetic} --- $16$-sample blocks (15 Gaussian $+$ one outlier at
$\kappa\!\cdot\!\rms$); bulk RMSE on the $15$ Gaussian entries:
\begin{center}\small
\begin{tabular}{@{}lcccccccc@{}}
\toprule
$\kappa$ & 2.0 & 2.5 & 3.0 & 4.0 & 6.0 & 8.0 & 14.0 & 20.0 \\
\midrule
GF4   & .102 & .103 & .105 & .110 & .130 & .154 & .232 & .309 \\
NVFP4 & .102 & .108 & .111 & .122 & .151 & .199 & .325 & .469 \\
ratio & 1.00 & 1.04 & 1.06 & 1.11 & 1.16 & 1.29 & 1.40 & 1.52 \\
\bottomrule
\end{tabular}
\end{center}
Tie at the Gaussian base $\kappa\!\approx\!2.5$; GF4's edge grows monotonically
with $\kappa$, exactly Theorem~\ref{thm:gf4}.

\emph{(B) Real} --- rotated OPT-125m \texttt{fc2} activation blocks have
$\kappa$: median $2.06$, $90$th pct $2.52$, max $3.73$ (rotation keeps $\kappa$
small). Binning by $\kappa$: $[2,3)$ ratio $1.07$, $[3,4)$ ratio $1.13$ --- a
modest but consistent GF4 win that grows with $\kappa$, as predicted.

\subsection{What is and is not claimed}

\begin{remark}[Scope and honest effect size]
The dominant lever is the \emph{scale rule} (RMS vs.\ absmax), not the codebook
shape (Lemma~\ref{lem:codebook}). Because rotation keeps $\kappa$ small
(median $2.06$), the \emph{per-block} granularity edge post-rotation is modest
($\sim\!7$--$13\%$ bulk RMSE); it is large only where $\kappa$ is large ---
pre-rotation, or on outlier-heavy blocks/layers. Two further GF4 advantages are
\emph{separate} from this granularity argument and compound with it: (i) GF4's
\emph{float} RMS scale lets residual multipass keep decaying, whereas NVFP4's
E4M3-quantized scale saturates ($\sim\!42$ dB), and (ii) the codebook match of
Lemma~\ref{lem:codebook}. The end-to-end perplexity ordering
(GF4 $<$ NVFP4 across the suite) is the sum of these three, each individually
small at post-rotation $\kappa$ but consistently signed.
\end{remark}

\end{document}
